\documentclass[12pt]{article}
\usepackage[utf8]{inputenc}
\usepackage[T1]{fontenc}
\usepackage{hyperref}
\title{Southeast Babar Language Resources for K-12 Curriculum Development}
\author{Compiled by Claude Sonnet 5 (claude-sonnet-5) for Andr\'as Kornai}
\date{}
\begin{document}
\maketitle

\section{Introduction}
Southeast Babar (ISO 639-3 vbb) is an Austronesian language of Babar Island, South Maluku, Indonesia, with about 4,500 speakers (2007). It belongs to the South Babar group of the Babar branch, with East Masela, West Masela and Serili. Search found only classification, word-list and phonology material; no educational resources.

\section{Primary Educational Resources}
No primer, textbook or school programme was found in search results.

\section{Linguistic and Reference Materials}
\begin{itemize}
\item Wikipedia, ``Southeast Babar language''. \url{https://en.wikipedia.org/wiki/Southeast_Babar_language}
\item Wikipedia, ``Babar languages'': classification of the branch. \url{https://en.wikipedia.org/wiki/Babar_languages}
\item Steinhauer, paper on East Nusantara language contact (PDF, SIL ICAL abstracts), mentions Babar phonology. \url{https://sil-philippines-languages.org/ical/abstracts/Steinhauer_East%20Nusantara%20lg%20contact.pdf}
\item Austronesian Basic Vocabulary Database, South-East Babar word list. \url{https://abvd.eva.mpg.de/austronesian/language.php?id=310}
\item ASJP word list for Southeast Babar. \url{https://asjp.clld.org/languages/SOUTHEAST_BABAR}
\item Verbix page with a short word list. \url{https://docs.verbix.com/Languages/BabarSoutheast}
\item Glottolog entry (sout2883). \url{https://glottolog.org/resource/languoid/id/sout2883.bigmap.html}
\end{itemize}

\section{Religious and Community Resources}
None found in search results.

\section{Digital and Audio Resources}
No audio, video or corpus resources were found in search results.

\section{Curriculum Development Context}
The language is endangered and is not known to be taught in schools. The available material is limited to word lists and linguistic classification. A curriculum would need new documentation with older speakers, and the full grammatical descriptions by Hein Steinhauer would need to be located through libraries; they were not found online in this search.

\end{document}